% !TeX program = xelatex
\documentclass[12pt]{article}
\usepackage[spanish, es-nodecimaldot]{babel}

\usepackage{fontspec}
\setmainfont{Times New Roman}

\usepackage[letterpaper,top=2.4cm,bottom=2.4cm,left=2.4cm,right=2.4cm]{geometry}

\usepackage{setspace}
% Renglon cerrado
\singlespacing
\usepackage{xurl}
\usepackage[hidelinks]{hyperref}

\usepackage{enumitem}

\usepackage{titlesec}

\setlength{\parindent}{0pt}

\setlength{\parskip}{6pt}

\titleformat{\section}
	{\normalfont\bfseries\fontsize{12}{14}\selectfont}
	{}{0pt}{}
\titlespacing*{\section}{0pt}{6pt}{0pt}

\author{Jenny Sofía Morales}

\begin{document}
	
	\begin{center}
		
		{\bfseries\fontsize{12}{14}\selectfont
			PROGRAMACIÓN REACTIVA Y SISTEMAS REACTIVOS
		}
		
		\vspace{0.4cm}
		
		{\itshape J. S. Morales López}
		
		\vspace{0.2cm}
		
		{\itshape 7690-08-6790 Universidad Mariano Gálvez}
		
		{\itshape Seminario de tecnologías de información}
		
		{\itshape
			\href{mailto:jmorales115@miumg.edu.gt}
			{jmorales115@miumg.edu.gt}
		}
		
		{\itshape 
			\href{https://github.com/sofiaMorales3805/ForoAcademico-Seminario}
			{Repositorio en GitHub - Sofia Morales}
		}
	\end{center}
	
	\vspace{0.5cm}
	\section*{Resumen}
	
	El presente artículo analiza los principios y las buenas prácticas necesarias para desarrollar aplicaciones reactivas capaces de responder oportunamente ante eventos, cambios de carga y fallos. A partir de la revisión de fuentes especializadas, se diferencia la programación reactiva, empleada para administrar flujos de datos asíncronos y no bloqueantes, de los sistemas reactivos, entendidos como un enfoque arquitectónico completo. Entre los principales resultados se identificaron cuatro propiedades: capacidad de respuesta, resiliencia, elasticidad y orientación a mensajes. También se reconocieron prácticas como documentar el flujo de eventos, aplicar contrapresión, reducir el estado compartido, aislar los fallos, incorporar métricas y ejecutar pruebas de carga y recuperación. Estas medidas favorecen el uso eficiente de los recursos y la continuidad del servicio; sin embargo, requieren una planificación cuidadosa, porque la asincronía puede dificultar la depuración y el monitoreo. Se concluye que una aplicación no se vuelve reactiva solamente por utilizar tareas asíncronas, sino por integrar principios de diseño, control del flujo, observabilidad y recuperación ante fallos.
	
	
	\noindent
	\textbf{Palabras clave:} programación reactiva, sistemas reactivos, resiliencia, elasticidad, mensajería asíncrona.
	
	
	\section*{Desarrollo del tema}
	
	La programación reactiva y los sistemas reactivos responden a necesidades
	relacionadas, pero no representan exactamente el mismo concepto. La
	programación reactiva se aplica principalmente al procesamiento de flujos de
	datos y eventos asíncronos, mientras que los sistemas reactivos incorporan
	estos principios dentro de un enfoque arquitectónico más amplio (Bonér, 2023).
	
	Para implementar este enfoque adecuadamente, las buenas prácticas deben
	considerarse desde la etapa de diseño. El flujo de eventos debe mostrar el
	origen de los datos, los componentes encargados de procesarlos y las respuestas
	generadas. Esto facilita la identificación de errores, retrasos y posibles
	puntos de saturación. Gillis y Nolle (2024) también recomiendan mantener
	actualizado el diagrama del flujo y revisar el acceso a las bases de datos para
	evitar latencia innecesaria.
	
	Una práctica especialmente importante es la contrapresión. Este mecanismo
	permite que el consumidor indique cuánto puede procesar y evita que un
	productor envíe más información de la que el sistema puede atender. De esta
	manera, se reduce el riesgo de utilizar recursos sin límite y de acumular
	eventos pendientes durante periodos de alta demanda (Bonér, 2023).
	
	También es necesario aislar los componentes para impedir que el fallo de uno
	afecte al sistema completo. El monitoreo mediante registros, métricas y
	alertas permite detectar comportamientos anormales y conocer el estado de la
	aplicación. Finalmente, las pruebas de carga, concurrencia y recuperación
	ayudan a comprobar si el sistema conserva su capacidad de respuesta ante
	fallos o variaciones en la cantidad de solicitudes.
	
	\begin{itemize}
	
	\item Diseñar y documentar el flujo de eventos.
	\item Emplear operaciones asíncronas y no bloqueantes.
	\item Aplicar contrapresión para evitar que un productor sature al consumidor.
	\item Reducir el estado compartido entre componentes.
	\item Aislar los fallos y establecer mecanismos de recuperación.
	\item Implementar registros, métricas y monitoreo.
	\item Realizar pruebas de carga, concurrencia y fallos.
	\item Evaluar si el enfoque reactivo realmente es necesario para la aplicación.
	
	\end{itemize}
	
	TechTarget recomienda diagramar los flujos de eventos, reducir componentes con estado y revisar el acceso a las bases de datos para evitar latencia (Gillis \& Nolle, 2024). Akka también explica que la programación reactiva es una técnica interna, mientras que un sistema reactivo representa un enfoque arquitectónico más amplio (Bonér, 2023).
	
	\section*{Programación reactiva}
	
	La programación reactiva permite gestionar flujos de información y responder
	a los cambios que se producen en los datos. Se apoya en operaciones asíncronas
	y no bloqueantes para que la aplicación pueda continuar atendiendo otros
	eventos mientras espera que finalice una operación, como una consulta a una
	base de datos o una solicitud a un servicio externo.
	
	Según Pérez (2026), la programación reactiva permite que los componentes
	respondan a modificaciones en los datos y procesen eventos mediante flujos
	asíncronos sin detener la ejecución de la aplicación. Sus principales
	principios son:
	
	\begin{itemize}
		\item Flujos de datos observables.
		\item Propagación de cambios.
		\item Manejo eficiente de eventos.
	\end{itemize}
	
	\section*{Sistemas reactivos}
	
	Los sistemas reactivos,son más flexibles, con bajo acoplamiento y escalables.	Son significativamente tolerantes a fallos y por ende cuando fallan responden de manera oportuna. Estos sistemas presentan cuatro características fundamentales según	(Bonér et al., 2014):
	
	\begin{itemize}
		
		\item \textbf{Responsivos:} mantienen tiempos de respuesta oportunos
		y consistentes para proporcionar un servicio de calidad.
		
		\item \textbf{Resilientes:} conservan su capacidad de respuesta ante
		fallos mediante el aislamiento, la contención y la recuperación de los
		componentes afectados.
		
		\item \textbf{Elásticos:} aumentan o reducen los recursos utilizados
		según los cambios en la carga de trabajo.
		
		\item \textbf{Orientados a mensajes:} utilizan mensajes asíncronos
		para reducir el acoplamiento entre componentes y controlar el flujo
		de información.
		
	\end{itemize}
	
	\section*{Observaciones y comentarios}
	
	El enfoque reactivo puede ser útil cuando una aplicación debe procesar
	eventos constantes, responder en tiempo real o adaptarse a cambios de
	carga. Sin embargo, su implementación también incrementa la complejidad
	del desarrollo, especialmente durante la depuración y el seguimiento de
	operaciones asíncronas. Por esta razón, no debería utilizarse únicamente
	porque sea una tecnología moderna. Antes de adoptarlo, es necesario
	analizar las necesidades, el volumen de eventos y los recursos disponibles.
	
	\section*{Conclusiones}
	
	\begin{enumerate}
		
		\item La programación reactiva facilita el procesamiento de eventos y
		flujos de datos asíncronos sin bloquear innecesariamente la ejecución
		de la aplicación.
			
		\item La contrapresión permite controlar la cantidad de información
		procesada y evita que un productor sature la capacidad del consumidor.
		
		\item La resiliencia, la elasticidad y el aislamiento de fallos permiten
		conservar la continuidad del servicio ante errores o cambios en la demanda.
		
		\item El enfoque reactivo puede utilizarse en aplicaciones de IoT, sistemas
		distribuidos y servicios en tiempo real; sin embargo, su implementación
		debe evaluarse según las necesidades y la complejidad del proyecto.
		
	\end{enumerate}
	
\section*{Bibliografía}

\begin{itemize}[
	leftmargin=1.27cm,
	label={},
	itemindent=-1.27cm,
	itemsep=6pt
	]
	
	\item Bonér, J. (2023, 24 de octubre).
	\textit{Reactive programming vs. reactive systems}.
	Akka.
	\url{https://akka.io/blog/reactive-programming-versus-reactive-systems}
	
	\item Bonér, J., Farley, D., Kuhn, R., \& Thompson, M.
	(2014, 16 de septiembre).
	\textit{El Manifiesto de Sistemas Reactivos}.
	\url{https://www.reactivemanifesto.org/es}
	
	\item Gillis, A. S., \& Nolle, T. (2024, 14 de junio).
	\textit{What is reactive programming?}
	TechTarget.
	\url{https://www.techtarget.com/searchapparchitecture/definition/reactive-programming}
	
	\item Pérez, A. (2026, 11 de mayo).
	\textit{Cómo aplicar la programación reactiva en aplicaciones web modernas}.
	OBS Business School.
	\url{https://www.obsbusiness.school/blog/como-aplicar-la-programacion-reactiva-en-aplicaciones-web-modernas-cp}
	\end{itemize}
	
\end{document}